\documentclass[12pt]{article}
\usepackage[utf8]{inputenc}
\usepackage[T1]{fontenc}
\usepackage[margin=1in]{geometry}
\usepackage{setspace}
\usepackage{xcolor}
\usepackage[colorlinks=true, urlcolor=blue!60!black, linkcolor=blue!60!black]{hyperref}
\setlength{\parindent}{2em}
\setlength{\parskip}{1.25em}
\renewcommand{\baselinestretch}{1.0}
\title{\textit{Mahmoud v. Taylor}, 606 U.S. 522 (2025)}
\author{Case Brief}
\date{}
\begin{document}
\maketitle

\section{Case Metadata}

\begin{description}
  \item[Case Name:] \textit{Mahmoud v. Taylor} (Tamer Mahmoud, et al., Petitioners, v. Thomas W. Taylor, Superintendent of Montgomery County Public Schools, et al., Respondents)
  \item[Citation:] 606 U.S. 522 (2025)
  \item[Docket No.:] No. 24-297
  \item[Court:] Supreme Court of the United States (on writ of certiorari to the U.S. Court of Appeals for the Fourth Circuit, 102 F.4th 191 (2024), affirming 688 F. Supp. 3d 265 (D. Md. 2023))
  \item[Argued:] April 22, 2025
  \item[Decided:] June 27, 2025
  \item[Vote:] 6--3 (reversed and remanded)
  \item[Majority Opinion Author:] Justice Samuel A. Alito, Jr. (joined by Chief Justice John G. Roberts, Jr., and Justices Clarence Thomas, Neil M. Gorsuch, Brett M. Kavanaugh, and Amy Coney Barrett)
  \item[Concurring Opinion:] Justice Clarence Thomas
  \item[Dissenting Opinion:] Justice Sonia Sotomayor (joined by Justices Elena Kagan and Ketanji Brown Jackson)
  \item[Constitutional Provision(s):] U.S. Const. amend. I (Free Exercise Clause), incorporated via amend. XIV
  \item[Chapter:] Religious Liberty
  \item[Primary Opinion Sources:] \href{https://www.law.cornell.edu/supremecourt/text/24-297}{Cornell LII} $|$ \href{https://supreme.justia.com/cases/federal/us/606/24-297/}{Justia} $|$ \href{https://www.supremecourt.gov/opinions/24pdf/24-297_4f14.pdf}{SupremeCourt.gov}
\end{description}

\section{Facts and Procedural History}

During the 2022--2023 school year, the Montgomery County Board of Education (Board) in Maryland---serving more than 160,000 students in one of the Nation's most religiously diverse counties---introduced five ``LGBTQ+-inclusive'' storybooks into its kindergarten through fifth-grade English Language Arts curriculum: \textit{Intersection Allies}, \textit{Prince \& Knight}, \textit{Love, Violet}, \textit{Born Ready: The True Story of a Boy Named Penelope}, and \textit{Uncle Bobby's Wedding} (\href{https://www.law.cornell.edu/supremecourt/text/24-297}{606 U.S. 522 (2025)}, slip op. at 2--6). Teachers were expected to use the books in instruction and received professional-development guidance encouraging them to ``[d]isrupt the either/or thinking'' of students and to correct comments questioning same-sex marriage or gender identity as ``hurtful'' (slip op. at 6--7).

Initially, consistent with its ``Guidelines for Respecting Religious Diversity,'' the Board notified parents when the storybooks would be taught and allowed religious opt-outs (slip op. at 8--9). In March 2023, however, the Board abruptly rescinded parental notice and opt-outs, citing classroom disruption, high absenteeism, and concerns about social stigma, even while continuing to allow opt-outs for Maryland's mandatory ``Family Life and Human Sexuality'' unit and noncurricular activities (slip op. at 9--11).

Three sets of religiously diverse parents---Tamer Mahmoud and Enas Barakat (Muslim), Jeff and Svitlana Roman (Roman Catholic and Ukrainian Orthodox), and Chris and Melissa Persak (Catholic)---along with the parents' association Kids First, sued in the U.S. District Court for the District of Maryland seeking a preliminary injunction restoring notice and opt-outs under the Free Exercise Clause (slip op. at 11--14). Both the Mahmouds and the Romans felt religiously compelled to withdraw their elementary-aged sons to private school at significant expense (slip op. at 11--12). The District Court denied relief (\textit{Mahmoud v. McKnight}, 688 F. Supp. 3d 265 (D. Md. 2023)), and a divided Fourth Circuit panel affirmed (102 F.4th 191 (4th Cir. 2024)), holding that plaintiffs failed to show coercion on a ``threadbare'' record (slip op. at 15--16). The Supreme Court granted certiorari and reversed.

\section{Issue(s) Presented}

Does a public school board unconstitutionally burden parents' Free Exercise rights under the First Amendment when it mandates elementary-school instruction using normative LGBTQ+-inclusive storybooks while withholding parental notice and religious opt-outs?

\section{Holding}

\textbf{Yes}; judgment \textbf{reversed and remanded} (6--3) (\href{https://www.law.cornell.edu/supremecourt/text/24-297}{606 U.S. 522 (2025)}, slip op. at 1--2, 40--41). Denying parents notice and religious opt-outs from normative elementary instruction on gender and sexuality ``substantially interferes with the religious development'' of their children under \href{https://supreme.justia.com/cases/federal/us/406/205/}{\textit{Wisconsin v. Yoder}, 406 U.S. 205, 218 (1972)}, triggers strict scrutiny regardless of facial neutrality or general applicability under \href{https://supreme.justia.com/cases/federal/us/494/872/}{\textit{Employment Division v. Smith}, 494 U.S. 872, 881 (1990)}, and fails narrow tailoring. Petitioners are entitled to a preliminary injunction requiring advance notice and opt-outs.

\section{Reasoning}

\subsection{Substantial Interference Under \textit{Wisconsin v. Yoder}}
Writing for the Court, Justice Alito anchored the analysis in parents' longstanding Free Exercise right ``to direct `the religious upbringing' of their children'' (\href{https://www.law.cornell.edu/supremecourt/text/24-297}{slip op. at 17--18}, quoting \href{https://supreme.justia.com/cases/federal/us/406/205/}{\textit{Yoder}, 406 U.S. at 213--214}). While \textit{West Virginia State Board of Education v. Barnette}, 319 U.S. 624 (1943), addressed direct compulsion to affirm a belief, \textit{Yoder} confirmed that the Free Exercise Clause also bars subtler governmental policies that place children in an environment ``hostile'' to their faith with ``pressure to conform,'' posing ``a very real threat of undermining'' parental religious instruction (slip op. at 19--21, quoting \textit{Yoder}, 406 U.S. at 211, 218).

Evaluating the context and the young age (5 to 11) of impressionable elementary students, the Court found the storybooks ``unmistakably normative'': \textit{Prince \& Knight} and \textit{Uncle Bobby's Wedding} present same-sex marriage as an occasion for universal celebration, while \textit{Intersection Allies} and \textit{Born Ready} endorse fluid gender identity and rebuke biological-sex skepticism (slip op. at 21--25). Paired with teacher guidance encouraging educators to ``disrupt'' student thinking and label contrary views ``hurtful,'' the curriculum went far beyond neutral ``exposure'' and posed an ``objective danger'' to religious upbringing (slip op. at 25--28).

\subsection{Rejection of the Board's Counterarguments}
The Court rejected four defenses: (1) \textit{Bowen v. Roy}, 476 U.S. 693 (1986), and \textit{Lyng v. Northwest Indian Cemetery Protective Ass'n}, 485 U.S. 439 (1988), govern the government's ``internal affairs,'' not direct, compulsory interactions with public school children (slip op. at 28--29); (2) \textit{Yoder} is not ``sui generis'' to the Amish but states a general Free Exercise principle (slip op. at 29--31); (3) plaintiffs need not wait for classroom harm to pursue a pre-enforcement First Amendment challenge when Board policy is undisputed (slip op. at 31--32); and (4) the availability of private school, homeschooling, or after-school religious instruction cannot cure an unconstitutional condition on the public benefit of free, compulsory public education (slip op. at 32--35).

\subsection{Bypassing \textit{Smith} and Applying Strict Scrutiny}
Although neutral, generally applicable laws normally receive rational-basis review under \href{https://supreme.justia.com/cases/federal/us/494/872/}{\textit{Smith}, 494 U.S. at 878--879}, the Court held that when a policy imposes a burden ``of the same character as that in \textit{Yoder},'' strict scrutiny applies automatically without needing to evaluate neutrality, general applicability, or ``hybrid rights'' (\href{https://www.law.cornell.edu/supremecourt/text/24-297}{slip op. at 35--37 \& n.14}). Applying strict scrutiny under \href{https://supreme.justia.com/cases/federal/us/593/19-123/}{\textit{Fulton v. City of Philadelphia}, 593 U.S. 522, 541 (2021)}, and \href{https://supreme.justia.com/cases/federal/us/508/520/}{\textit{Church of the Lukumi Babalu Aye, Inc. v. City of Hialeah}, 508 U.S. 520, 546 (1993)}, the Court held the no-opt-out policy failed narrow tailoring: the Board already accommodated opt-outs for sex education and pull-out programs for multilingual and special-education students, its administrative burdens were ``self-inflicted'' by weaving the books throughout ELA, and ``[t]here is no \textit{de maximis} exception to the Free Exercise Clause'' (slip op. at 37--40). Finding irreparable First Amendment injury, the Court ordered preliminary injunctive relief (slip op. at 40--41).

\section{Concurring and Dissenting Opinions}

\begin{itemize}
  \item \textbf{Justice Thomas (Concurring):} Joining in full, Justice Thomas argued that under \textit{Yoder}'s historical inquiry, elementary sexuality and gender-identity instruction lacks any historical pedigree and is inessential to civic duties. Analogizing the Board's equity rationale to the assimilationist conformity struck down in \textit{Pierce v. Society of Sisters}, 268 U.S. 510 (1925), he warned that allowing districts to defeat opt-outs by weaving controversial themes into core classes would ``give schools a playbook for evading the First Amendment'' (\href{https://www.law.cornell.edu/supremecourt/text/24-297}{Thomas, J., concurring, slip op. at 1--12}).
  \item \textbf{Justice Sotomayor (Dissenting, joined by Justices Kagan and Jackson):} Justice Sotomayor argued that the Free Exercise Clause bars governmental coercion, not mere ``exposure'' to ideas conflicting with parental faith, and that \textit{Lyng} and \textit{Bowen} foreclosed the majority's ``threat of undermining'' test. She contended that the Board's guidance promoted civility rather than coercion, that bypassing \textit{Smith} without a hybrid due-process analysis distorted precedent, and that mandating individualized notice and opt-outs for ``subtle'' normative themes would inflict administrative ``chaos'' and chill public school curricula nationwide (\href{https://www.law.cornell.edu/supremecourt/text/24-297}{Sotomayor, J., dissenting, slip op. at 1--38}).
\end{itemize}

\section{Key Quotations}

\begin{enumerate}
  \item \textbf{Core Holding on Parental Religious Upbringing} (\href{https://www.law.cornell.edu/supremecourt/text/24-297}{606 U.S. 522 (2025)}, slip op. at 1--2; \href{https://supreme.justia.com/cases/federal/us/606/24-297/}{Justia}):
  \begin{quote}
    ``A government burdens the religious exercise of parents when it requires them to submit their children to instruction that poses `a very real threat of undermining' the religious beliefs and practices that the parents wish to instill. \textit{Wisconsin v. Yoder}, 406 U. S. 205, 218 (1972). And a government cannot condition the benefit of free public education on parents' acceptance of such instruction.''
  \end{quote}

  \item \textbf{Automatic Strict Scrutiny for \textit{Yoder}-Style Burdens} (\href{https://www.law.cornell.edu/supremecourt/text/24-297}{slip op. at 36}; \href{https://www.supremecourt.gov/opinions/24pdf/24-297_4f14.pdf}{SupremeCourt.gov}):
  \begin{quote}
    ``When the burden imposed is of the same character as that imposed in \textit{Yoder}, we need not ask whether the law at issue is neutral or generally applicable before proceeding to strict scrutiny. That much is clear from our decisions in \textit{Yoder} and \textit{Smith}.''
  \end{quote}

  \item \textbf{No ``De Maximis'' Exception to Free Exercise} (\href{https://www.law.cornell.edu/supremecourt/text/24-297}{slip op. at 39}; \href{https://www.supremecourt.gov/opinions/24pdf/24-297_4f14.pdf}{SupremeCourt.gov}):
  \begin{quote}
    ``When it comes to instruction that would burden the religious exercise of parents, the Board cannot escape its obligations under the Free Exercise Clause by crafting a curriculum that is so burdensome that a substantial number of parents elect to opt out. There is no de maximis exception to the Free Exercise Clause.''
  \end{quote}

  \item \textbf{Dissenting Warning Against an Exposure-Based Veto} (\href{https://www.law.cornell.edu/supremecourt/text/24-297}{Sotomayor, J., dissenting, slip op. at 1--2}; \href{https://supreme.justia.com/cases/federal/us/606/24-297/}{Justia}):
  \begin{quote}
    ``Casting aside longstanding precedent, the Court invents a constitutional right to avoid exposure to `subtle' themes `contrary to the religious principles' that parents wish to instill in their children. [Internal citation omitted.] Exposing students to the `message' that LGBTQ people exist, and that their loved ones may celebrate their marriages and life events, the majority says, is enough to trigger the most demanding form of judicial scrutiny.''
  \end{quote}
\end{enumerate}

\section{Doctrinal Significance and Modern Legacy (Through 2026)}

\textit{Mahmoud} dramatically revitalized \href{https://supreme.justia.com/cases/federal/us/406/205/}{\textit{Yoder}} and carved a major categorical bypass around \href{https://supreme.justia.com/cases/federal/us/494/872/}{\textit{Employment Division v. Smith}}: without overruling \textit{Smith}, the Court established that any public-school policy imposing a \textit{Yoder}-type threat to parental religious formation triggers strict scrutiny automatically. It also rejected decades of lower-court consensus (such as \textit{Mozert v. Hawkins County Board of Education}, 827 F.2d 1058 (6th Cir. 1987)) that had treated non-coercive classroom exposure as outside the Free Exercise Clause.

Within a year, the Supreme Court expanded \textit{Mahmoud}'s reach beyond classroom storybooks. On March 2, 2026, in \href{https://www.law.cornell.edu/supremecourt/text/25A810}{\textit{Mirabelli v. Bonta}, 607 U.S. 492 (2026) (per curiam)} (\href{https://www.supremecourt.gov/opinions/25pdf/607us2r17_bqmc.pdf}{SupremeCourt.gov}), the Court vacated a Ninth Circuit stay and reinstated an injunction blocking California school policies that barred staff from disclosing a student's social gender transition to parents without student consent. Relying squarely on \textit{Mahmoud}, the Court held that concealing a child's gender transition imposes an even ``greater'' intrusion on parents' Free Exercise and Fourteenth Amendment substantive due process rights and likely fails \textit{Mahmoud}'s strict scrutiny. Across 2025--2026, public school districts nationwide---including Montgomery County Public Schools---revised parental-notice and opt-out policies to avoid strict-scrutiny litigation (\href{https://www.k12dive.com/news/what-the-supreme-courts-parental-opt-out-ruling-means-for-schools/753608/}{K-12 Dive}).

\section{Pedagogical Critiques and Exam Hypotheticals}

\subsection{Pedagogical Critiques}
\begin{enumerate}
  \item \textbf{Erosion of \textit{Smith} Without Overruling It:} By treating \textit{Yoder} as a standalone trigger for strict scrutiny without requiring a ``hybrid'' due-process claim, \textit{Mahmoud} narrows \textit{Smith} substantially while leaving lower courts uncertain which other Free Exercise burdens qualify as ``\textit{Yoder}-like.''
  \item \textbf{Line-Drawing and Administrability:} Distinguishing impermissibly ``normative'' elementary instruction from neutral ``exposure'' forces federal judges to parse children's books and teacher scripts, raising Justice Sotomayor's concern that resource-strapped districts may self-censor curricula rather than manage individualized opt-outs.
\end{enumerate}

\subsection{Exam Hypotheticals}
\textbf{1. Applying \textit{Mahmoud} (Parents Prevail):} A public school district introduces a mandatory second-grade civics workbook titled \textit{No Gods Needed} that portrays religious belief as outdated superstition and directs teachers to respond to students who cite Scripture by saying, ``Those stories are just myths.'' No notice or opt-out is allowed, though the district allows opt-outs for health class. Under \textit{Mahmoud}, normative instruction hostile to parents' faith aimed at impressionable seven-year-olds triggers automatic strict scrutiny under \textit{Yoder} and fails narrow tailoring given existing opt-out mechanisms.

\textbf{2. ``Opposite Side'' Contrast (School District Prevails):} A public high school requires eleventh-grade biology students to study evolutionary theory from a standard, neutral science textbook that neither disparages religion nor instructs teachers to rebuke religious students, and explicitly permits students to hold personal religious views. A parent sues for a religious opt-out under \textit{Mahmoud}. \textbf{Analysis:} The district has strong grounds to prevail. Unlike \textit{Mahmoud}'s five-to-eleven-year-olds, high school juniors are far less impressionable (\href{https://www.law.cornell.edu/supremecourt/text/24-297}{slip op. at 21--22}); the instruction is neutrally presented academic science rather than normative moral advocacy with teacher-led ``disruption'' of religious views; and core high-school science standards advance a compelling educational interest where piecemeal opt-outs are far harder to administer without undermining graduation requirements.

\end{document}
